\documentclass{article}

\usepackage{biblatex}
\usepackage{parskip}

\addbibresource{bib/magic.bib}
\addbibresource{bib/local.bib}

%%%%%%%%%%%%%%%%%%%%%%%%%%%%%%%%%%%%%%%
% Variant-dependent Config and Macros %
%%%%%%%%%%%%%%%%%%%%%%%%%%%%%%%%%%%%%%%

% Fancy \byvariant macro
\ExplSyntaxOn
\msg_new:nnn { byvariant } { unhandled } { Unhandled~variant:~#1 }
\NewExpandableDocumentCommand{\byvariant}{o+m}
  {
    \str_case_e:nnF { \variant } {#2}
      {
        \tl_if_novalue:nTF {#1}
          { \exp_args:NnnV \msg_expandable_error:nnn { byvariant } { unhandled } \variant }
          {#1}
      }
  }
\ExplSyntaxOff

\newcommand{\prompt}[1]{%
  \byvariant{
    {draft}{{\color{RoyalBlue}#1}}
    {final}{}
  }%
}

\begin{document}

\section{Section}

A citation from DOI~\cite{DOI:10.1145/3360598}, from
MLR~\cite{MLR:v267/park25l} and one from the local bib~\cite{hopcroft-karp}.

You can also use jinja for citation aliases, defined in
\texttt{bib/aliases.toml}~\cite{<<Duet>>}.

Here is a jinja template that pulls data from \texttt{data/data.toml}...

\begin{tabular}{|c|}\hline
<@ for x in xs @>
<< x >> \\\hline
<@ endfor @>
\end{tabular}

\printbibliography

\end{document}
